\documentclass[11pt,oneside,openany]{book}

% =============================================================================
% Configuración general
% =============================================================================
\usepackage[utf8]{inputenc}
\usepackage[T1]{fontenc}
\usepackage[provide=*,spanish]{babel}
\usepackage{lmodern}
\usepackage[a4paper,top=2.35cm,bottom=2.45cm,left=2.45cm,right=2.20cm,
            headheight=15pt,footskip=1.25cm]{geometry}
\linespread{1.06}
\setlength{\parskip}{0.55em}
\setlength{\parindent}{0pt}

% =============================================================================
% Matemática, tablas, figuras y recursos visuales
% =============================================================================
\usepackage{amsmath,amssymb}
\usepackage{graphicx}
\graphicspath{{assets/}}
\usepackage{xcolor}
\usepackage{booktabs}
\usepackage{array}
\usepackage{tabularx}
\usepackage{longtable}
\usepackage{float}
\usepackage{hyperref}

\definecolor{AzulProfundo}{HTML}{102A43}
\definecolor{AzulMedio}{HTML}{176B87}
\definecolor{Turquesa}{HTML}{2A9D8F}
\definecolor{TurquesaOscuro}{HTML}{207A70}
\definecolor{Arena}{HTML}{F4E9D8}
\definecolor{ArenaClara}{HTML}{FBF6EE}
\definecolor{CelesteClaro}{HTML}{EAF4F8}
\definecolor{VerdeClaro}{HTML}{EAF7F3}
\definecolor{AmarilloClaro}{HTML}{FFF6DF}
\definecolor{Naranja}{HTML}{E9A23B}
\definecolor{GrisTexto}{HTML}{3F4D5A}
\definecolor{GrisClaro}{HTML}{EEF2F5}
\definecolor{RojoSuave}{HTML}{C95F5F}

\hypersetup{
  colorlinks=true,
  linkcolor=AzulMedio,
  citecolor=TurquesaOscuro,
  urlcolor=AzulMedio,
  pdftitle={Proyecto 2 - MAP: Charting Student Math Misunderstandings},
  pdfauthor={Jorge Gabriel Palacios Sales; Pablo Daniel Barillas Moreno; Roberto Emiliano Otoniel Camposeco Torres; Abby Sofia Donis Agreda},
  pdfsubject={Análisis exploratorio de explicaciones matemáticas estudiantiles mediante NLP}
}

\makeatletter
\def\@makechapterhead#1{%
  \vspace*{18pt}%
  {\parindent\z@\raggedright\normalfont
    \ifnum\c@secnumdepth>\m@ne
      \sffamily\bfseries\color{Turquesa}\fontsize{36}{40}\selectfont
      \thechapter\par\nobreak\vskip 8pt
    \fi
    \sffamily\bfseries\color{AzulProfundo}\Huge #1\par\nobreak
    \vskip 7pt\color{Turquesa}\hrule height 1.1pt\vskip 28pt}}
\def\@makeschapterhead#1{%
  \vspace*{18pt}%
  {\parindent\z@\raggedright\normalfont
    \sffamily\bfseries\color{AzulProfundo}\Huge #1\par\nobreak
    \vskip 7pt\color{Turquesa}\hrule height 1.1pt\vskip 28pt}}
\renewcommand\section{\@startsection{section}{1}{\z@}%
  {-3.5ex \@plus -1ex \@minus -.2ex}{2.3ex \@plus .2ex}%
  {\normalfont\Large\sffamily\bfseries\color{AzulMedio}}}
\renewcommand\subsection{\@startsection{subsection}{2}{\z@}%
  {-3.25ex\@plus -1ex \@minus -.2ex}{1.5ex \@plus .2ex}%
  {\normalfont\large\sffamily\bfseries\color{TurquesaOscuro}}}
\renewcommand\subsubsection{\@startsection{subsubsection}{3}{\z@}%
  {-3.25ex\@plus -1ex \@minus -.2ex}{1.5ex \@plus .2ex}%
  {\normalfont\normalsize\sffamily\bfseries\color{AzulProfundo}}}
\makeatother

% =============================================================================
% Componentes reutilizables
% =============================================================================
\newsavebox{\infobox}
\newenvironment{hallazgo}[1][]{%
  \begin{lrbox}{\infobox}\begin{minipage}{0.92\textwidth}
  \color{GrisTexto}\sffamily\textbf{#1}\par\smallskip\rmfamily}
  {\end{minipage}\end{lrbox}\par\medskip\noindent
   \colorbox{CelesteClaro}{\usebox{\infobox}}\par\medskip}
\newenvironment{decision}[1][]{%
  \begin{lrbox}{\infobox}\begin{minipage}{0.92\textwidth}
  \color{GrisTexto}\sffamily\textbf{#1}\par\smallskip\rmfamily}
  {\end{minipage}\end{lrbox}\par\medskip\noindent
   \colorbox{VerdeClaro}{\usebox{\infobox}}\par\medskip}
\newenvironment{advertencia}[1][]{%
  \begin{lrbox}{\infobox}\begin{minipage}{0.92\textwidth}
  \color{GrisTexto}\sffamily\textbf{#1}\par\smallskip\rmfamily}
  {\end{minipage}\end{lrbox}\par\medskip\noindent
   \colorbox{AmarilloClaro}{\usebox{\infobox}}\par\medskip}
\newenvironment{ficha}[1][]{%
  \begin{lrbox}{\infobox}\begin{minipage}{0.90\textwidth}
  \color{AzulProfundo}\sffamily\textbf{#1}\par\smallskip\rmfamily}
  {\end{minipage}\end{lrbox}\par\medskip\noindent
   \fcolorbox{Arena}{ArenaClara}{\usebox{\infobox}}\par\medskip}

\newcommand{\metrica}[2]{%
  \begin{minipage}[t]{0.48\linewidth}
    {\sffamily\scriptsize\color{GrisTexto}#1}\par
    {\sffamily\bfseries\color{AzulProfundo}#2}
  \end{minipage}%
}
\newcommand{\figura}[4][0.97\textwidth]{%
  \begin{figure}[htbp]
    \centering
    \includegraphics[width=#1]{#2}
    \caption{#3}
    \label{#4}
  \end{figure}
}
\newcolumntype{Y}{>{\raggedright\arraybackslash}X}
\newcolumntype{C}[1]{>{\centering\arraybackslash}p{#1}}
\newcolumntype{R}[1]{>{\raggedleft\arraybackslash}p{#1}}
\newcommand{\codigo}[1]{\texttt{#1}}

\begin{document}

% =============================================================================
% Portada
% =============================================================================
\begin{titlepage}
\pagecolor{AzulProfundo}

\color{white}
\vspace*{0.5cm}
{\sffamily\bfseries\Large Universidad del Valle de Guatemala\par}
\vspace{0.2cm}
{\sffamily Facultad de Ingeniería\par}
{\sffamily Departamento de Ciencias de la Computación\par}

\vfill

{\sffamily\bfseries\color{Turquesa}\Large Proyecto 2 -- Análisis Exploratorio\par}
\vspace{0.55cm}
{\sffamily\bfseries\fontsize{32}{37}\selectfont
MAP: Charting Student\par
Math Misunderstandings\par}
\vspace{0.55cm}
{\Large\color{Arena}
Caracterización de explicaciones matemáticas estudiantiles mediante
Procesamiento del Lenguaje Natural\par}

\vspace{1.05cm}
\noindent\fcolorbox{white!40!AzulProfundo}{AzulProfundo}{%
\begin{minipage}{0.94\textwidth}
\vspace{3mm}
\color{white}
\begin{tabularx}{\textwidth}{@{}>{\sffamily\bfseries}p{3.4cm}Y@{}}
Curso: & CC3084 -- Data Science\\[1.5mm]
Docente: & Lynette García Pérez\\[1.5mm]
Grupo y sección: & Grupo 1, sección 10\\[1.5mm]
Integrantes: & Jorge Gabriel Palacios Sales\\
& Pablo Daniel Barillas Moreno\\
& Roberto Emiliano Otoniel Camposeco Torres\\
& Abby Sofia Donis Agreda\\[1.5mm]
Tema: & Procesamiento del Lenguaje Natural (NLP)\\
\end{tabularx}
\vspace{3mm}
\end{minipage}}

\vfill
{\sffamily\large Guatemala, agosto de 2026\par}
\end{titlepage}
\nopagecolor

\frontmatter

\chapter*{Ficha del informe}
\addcontentsline{toc}{chapter}{Ficha del informe}

\begin{hallazgo}[Propósito]
Este informe presenta el análisis exploratorio del reto 6, \emph{MAP --
Charting Student Math Misunderstandings}. Su propósito es caracterizar la
estructura, calidad, desbalance y contenido lingüístico de las explicaciones
escritas por estudiantes, y traducir la evidencia en decisiones justificadas
de preprocesamiento, validación y modelado.
\end{hallazgo}

\begin{table}[H]
\centering
\caption{Identificación y alcance del documento.}
\begin{tabularx}{\textwidth}{@{}p{4.15cm}Y@{}}
\toprule
\textbf{Elemento} & \textbf{Descripción}\\
\midrule
Curso & CC3084 -- Data Science, sección 10\\
Período académico & Segundo semestre de 2026\\
Reto & 6: MAP -- Charting Student Math Misunderstandings\\
Área & Procesamiento del Lenguaje Natural\\
Unidad de análisis & Respuesta de un estudiante a una pregunta matemática,
acompañada de una explicación libre\\
Base principal & \codigo{train.csv}: 36,696 observaciones y 7 variables\\
Objetivo de la competencia & Predecir la etiqueta compuesta
\codigo{Category:Misconception}\\
Alcance del informe & Exploración, limpieza, interpretación y próximos pasos;
no incluye entrenamiento de un modelo final\\
Fuente de evidencia & Contenido, tablas, salidas y siete figuras almacenadas
en el notebook del proyecto\\
\bottomrule
\end{tabularx}
\end{table}

\begin{advertencia}[Trazabilidad del análisis]
Para elaborar este documento no se ejecutaron celdas del notebook. Las cifras
y figuras proceden de sus salidas ya almacenadas. Cuando una afirmación textual
del notebook difería de una tabla, una gráfica o una salida numérica guardada,
se priorizó la evidencia numérica y la discrepancia se documentó como
limitación de reproducibilidad.
\end{advertencia}

\chapter*{Resumen}
\addcontentsline{toc}{chapter}{Resumen}

Se analizaron 36,696 respuestas estudiantiles distribuidas entre 15 preguntas
matemáticas. Cada registro contiene el enunciado, la opción elegida y una
explicación breve en lenguaje natural; en entrenamiento también incluye la
categoría del razonamiento y, cuando aplica, la concepción errónea. El espacio
de salida contiene seis valores de \codigo{Category}, 35 concepciones erróneas
y 65 etiquetas compuestas.

La investigación previa muestra que las 36,696 filas locales son la porción
etiquetada de un benchmark de 52,463 explicaciones. El conjunto completo fue
construido con 38,095 explicaciones originales y 14,368 paráfrasis sintéticas
(27\%), calificadas con rúbricas por ítem y controles de acuerdo entre
evaluadores (Rittle-Johnson et al., 2025). Esta procedencia obliga a interpretar
la ortografía, las minúsculas y las repeticiones como una mezcla de producción
estudiantil y aumento controlado.

El conjunto está completo en sus variables obligatorias. Los 26,836 valores
nulos de \codigo{Misconception} (73.1\%) son ausencias estructurales: aparecen
exactamente cuando el razonamiento no fue clasificado como concepción errónea.
Por ello se codifican como \codigo{NA} y no se imputan. No hay filas idénticas
si se conserva el identificador único; sin embargo, existen 736 duplicados al
ignorar \codigo{row\_id} y 970 repeticiones adicionales de texto. La ausencia
de un identificador de estudiante impide demostrar independencia, de modo que
estas observaciones se conservan, pero se reconoce su posible efecto sobre una
partición aleatoria.

La variable objetivo presenta una cola larga severa: 32 de las 65 clases
compuestas tienen menos de 50 observaciones y 24 tienen menos de 10, mientras
\codigo{True\_Correct:NA} concentra 40.3\% del total. El texto también impone
retos particulares: la explicación mediana tiene 13 palabras, 62.8\% comienza
en minúscula, 78.3\% contiene dígitos y 27.2\% usa una barra de fracción. Por
consiguiente, eliminar números o símbolos matemáticos destruiría información
central.

El hallazgo metodológico más importante es la dependencia entre etiquetas y
preguntas. Treinta y dos de las 35 concepciones erróneas aparecen en una sola
pregunta; además, 58 de los 60 pares \codigo{(QuestionId, MC\_Answer)}
determinan la veracidad sin necesidad de leer el texto. Esto favorece un modelo
jerárquico, pero también crea riesgo de memorización y fuga contextual. Se
recomienda reportar dos esquemas de validación: uno estratificado para preguntas
conocidas y otro agrupado por pregunta para medir transferencia a preguntas
nuevas.

\begin{decision}[Conclusión ejecutiva]
El reto no es una clasificación plana homogénea. La veracidad de la opción es
casi determinista a partir de la pregunta y la respuesta; la dificultad real
consiste en interpretar la explicación y distinguir razonamiento correcto,
concepción errónea o evidencia insuficiente. El diseño de validación es tan
importante como el algoritmo y debe hacer visible la diferencia entre memorizar
una pregunta y generalizar el razonamiento matemático.
\end{decision}

\tableofcontents
\listoffigures
\listoftables

\mainmatter

% =============================================================================
\chapter{Introducción}

\section{Situación problemática}

Una respuesta matemática incorrecta revela que existe un problema, pero no
explica su causa. Para intervenir pedagógicamente, el docente necesita conocer
el razonamiento que condujo al resultado. Una \emph{misconception} o concepción
errónea es un patrón de pensamiento incorrecto y relativamente estable; no es
equivalente a un simple descuido aritmético ni a la ausencia de explicación.

Dos estudiantes pueden seleccionar la misma opción equivocada por causas
distintas: uno puede invertir numerador y denominador y otro contar de forma
incorrecta las partes de una figura. Aunque el resultado observable sea el
mismo, las intervenciones educativas deberían ser diferentes. También puede
ocurrir que un estudiante acierte por casualidad con un razonamiento erróneo, o
que razone correctamente y cometa un error final de cálculo. Una calificación
basada únicamente en la opción marcada no detecta estas diferencias.

La dificultad práctica es la escala. Leer y diagnosticar manualmente cada
explicación después de una evaluación demanda tiempo especializado. El reto,
organizado en torno a datos de la plataforma educativa Eedi, propone automatizar
una primera caracterización para apoyar retroalimentación oportuna. El sistema
no sustituye el juicio docente: busca priorizar y organizar evidencia escrita.

\begin{hallazgo}[Relevancia educativa]
Los 403 casos \codigo{True\_Misconception} son especialmente importantes: la
opción marcada es correcta, pero la explicación revela una concepción errónea.
Un examen que solo califique el resultado los consideraría éxitos y perdería la
oportunidad de intervenir.
\end{hallazgo}

\section{Problema científico}

\begin{ficha}[Pregunta de investigación]
¿Es posible predecir automáticamente, a partir de la explicación en lenguaje
natural escrita por un estudiante, el enunciado y la opción seleccionada, si su
razonamiento es correcto y, cuando no lo es, qué concepción errónea específica
lo explica?
\end{ficha}

Formalmente, se trata de clasificación multiclase sobre texto corto y contenido
matemático. Para cada observación (i), los predictores son:

\begin{equation}
X_i = \{Q_i,\; T_i,\; A_i,\; E_i\},
\end{equation}

donde $Q_i$ es el identificador de la pregunta, $T_i$ su enunciado,
$A_i$ la opción seleccionada y $E_i$ la explicación libre. La salida es la
etiqueta compuesta:

\begin{equation}
Y_i = \text{Category}_i : \text{Misconception}_i.
\end{equation}

Las dificultades anticipadas son texto corto y ruidoso, errores ortográficos,
notación matemática informal, desbalance severo, clases con soporte mínimo y
dependencia entre la pregunta y la etiqueta.

\section{Objetivos}

\subsection{Objetivo general}

Caracterizar mediante análisis exploratorio la estructura, calidad y principales
dificultades del conjunto de datos MAP, de modo que las decisiones de
preprocesamiento, validación y diseño del modelo de clasificación queden
fundamentadas en evidencia empírica.

\subsection{Objetivos específicos}

\begin{enumerate}
  \item Cuantificar el desbalance de \codigo{Category},
        \codigo{Misconception} y la etiqueta compuesta, identificando las clases
        con menos de 50 y menos de 10 observaciones.
  \item Evaluar la calidad y el poder discriminante del texto mediante longitud,
        densidad numérica, diversidad léxica y símbolos matemáticos, además de
        detectar valores atípicos y ausencias.
  \item Analizar la relación entre pregunta, opción seleccionada, veracidad,
        tipo de razonamiento y concepción errónea para identificar dependencias
        que condicionen la validación.
  \item Proponer decisiones de limpieza y próximos pasos de modelado que
        conserven la información matemática y reduzcan el riesgo de fuga.
\end{enumerate}

\section{Alcance}

El informe cubre investigación conceptual de NLP, descripción de datos,
limpieza, estadística descriptiva, frecuencias, proporciones, cruces,
correlaciones, atípicos, interpretación léxica, conclusiones y próximos pasos.
No entrena ni compara modelos finales. El \codigo{test.csv} local tiene solo
tres filas de ejemplo; cualquier evaluación futura deberá construirse con una
partición controlada de entrenamiento.

% =============================================================================
\chapter{Investigación previa y antecedentes}

\section{Por qué estudiar concepciones erróneas}

Una concepción errónea no es simplemente una respuesta incorrecta. Es una
forma sistemática de interpretar o aplicar conocimiento matemático que puede
producir el mismo patrón de error en situaciones relacionadas. Esta diferencia
es pedagógicamente central: una equivocación de cálculo puede requerir práctica
o revisión, mientras que una regla conceptual mal generalizada exige hacer
visible el razonamiento y reconstruir su significado.

La literatura reciente sobre diagnóstico asistido por inteligencia artificial
coincide en que las explicaciones abiertas contienen información que la opción
marcada no ofrece. En un benchmark independiente de álgebra de nivel medio,
Otero, Druga y Lan (2025) mostraron que restringir los candidatos por tema y
añadir revisión de educadores mejoró notablemente la identificación de
concepciones erróneas. Este resultado respalda dos decisiones del presente
proyecto: incorporar el contexto de la pregunta y conservar supervisión humana
en cualquier uso pedagógico.

Los aciertos con razonamiento defectuoso merecen atención particular. Un
trabajo reciente presentado en el taller PEAF estudió este fenómeno con
respuestas de Eedi y advirtió que, aun cuando modelos de razonamiento mejoran la
detección, los falsos positivos pueden superar ampliamente los casos reales;
por ello propuso un flujo de detectar, verificar y escalar, en vez de generar
diagnósticos automáticos definitivos (Imran \& Bulathwela, 2026). Esta evidencia
se considera emergente porque corresponde a un trabajo de taller, pero coincide
con la necesidad de interpretar con cautela la clase
\codigo{True\_Misconception}.

\section{El reto MAP}

La competencia \emph{MAP -- Charting Student Math Misunderstandings} fue
organizada por Vanderbilt University y The Learning Agency en Kaggle. Su
objetivo oficial fue construir modelos de NLP capaces de proponer concepciones
erróneas candidatas a partir de explicaciones abiertas, para facilitar que los
docentes identifiquen razonamientos incorrectos que son costosos de etiquetar
manualmente (King et al., 2025). La competencia inició el 10 de julio de 2025 y
cerró el 15 de octubre de 2025.

El reto no se limita a predecir si la opción es correcta. La salida integra:

\begin{enumerate}
  \item la veracidad de la opción elegida;
  \item la calidad de la explicación: correcta, concepción errónea o ninguna;
  \item el nombre de la concepción errónea cuando existe.
\end{enumerate}

La meta declarada es que las concepciones aprendidas puedan generalizar entre
problemas diferentes, no solo que se reproduzca el inventario de etiquetas de
una pregunta conocida (King et al., 2025). Esta aspiración vuelve especialmente
relevante el hallazgo exploratorio de que 32 de las 35 concepciones presentes
en el entrenamiento local aparecen en una sola pregunta.

El estudio retrospectivo de la entidad organizadora reportó más de 1,850
equipos, cerca de 40,000 envíos y resultados MAP@3 superiores a 0.94 entre las
soluciones principales. También destacó que la taxonomía fue desarrollada por
expertos de contenido para diferenciar corrección, errores comunes y
concepciones específicas (The Learning Agency, 2025). Un desempeño alto en la
competencia demuestra viabilidad técnica dentro de ese benchmark, pero no
garantiza por sí solo validez externa en nuevos currículos, idiomas o
poblaciones.

\section{Procedencia y construcción del benchmark}

Rittle-Johnson et al. (2025) documentaron académicamente la construcción del
conjunto MAP. El benchmark completo contiene 52,463 explicaciones sobre 15
preguntas de matemática de nivel medio, principalmente números racionales,
ecuaciones y funciones. Los ítems provienen de Eedi; cada uno contiene una
opción correcta y tres distractores diseñados, cuando es posible, para capturar
errores conocidos. A la selección de opción múltiple le sigue una solicitud de
explicación abierta.

\begin{table}[htbp]
\centering
\caption{Procedencia del benchmark MAP según su artículo metodológico.}
\label{tab:procedencia-benchmark}
\begin{tabularx}{\textwidth}{@{}p{4.5cm}Y@{}}
\toprule
\textbf{Elemento} & \textbf{Descripción documentada}\\
\midrule
Plataforma de origen & Eedi, plataforma de aprendizaje de matemáticas del Reino
Unido.\\
Cobertura & 15 preguntas de nivel medio, con énfasis en números racionales,
ecuaciones y funciones.\\
Tamaño completo & 52,463 explicaciones; el \codigo{train.csv} local contiene
36,696, es decir, solo la porción etiquetada disponible para entrenamiento.\\
Explicaciones originales & 38,095 registros.\\
Explicaciones aumentadas & 14,368 paráfrasis generadas con Coedit-XL, equivalentes
al 27\% del benchmark completo.\\
Anotación & Rúbricas específicas por ítem, desarrolladas por tres especialistas
en cognición matemática y concepciones erróneas.\\
Control de calidad & 24\% de explicaciones codificadas por dos evaluadores;
\(\kappa\) de Cohen entre 0.70 y 0.90.\\
Privacidad & Explicaciones revisadas para retirar información identificable; no
se distribuyeron datos demográficos.\\
\bottomrule
\end{tabularx}
\end{table}

La generación sintética buscó aumentar el soporte de concepciones poco
frecuentes y facilitar la capacitación de evaluadores. Se permitió como máximo
una paráfrasis por explicación auténtica. Para imitar la escritura estudiantil,
la mitad de los textos aumentados recibió un error ortográfico aleatorio y la
mitad fue convertida completamente a minúsculas (Rittle-Johnson et al., 2025).

\begin{advertencia}[Implicación para el análisis de calidad]
No todo el ruido observado puede atribuirse directamente a estudiantes. Parte
de las minúsculas, ortografía y repetición léxica fue introducida de manera
controlada en las explicaciones sintéticas. Como los archivos de competencia no
incluyen un indicador de procedencia, el análisis local no puede separar textos
originales de aumentados ni medir su distribución por clase.
\end{advertencia}

La anotación también impone límites. Cada explicación recibe un solo código,
aunque un razonamiento podría exhibir más de un error. El artículo reporta que
27\% del benchmark completo contiene evidencia de una posible concepción
errónea y que su frecuencia cambia considerablemente entre ítems
(Rittle-Johnson et al., 2025). Esto coincide con la cola larga y la dependencia
por pregunta encontradas en el EDA.

\section{Técnicas investigadas para detectar patrones escritos}

La investigación previa sugiere una progresión metodológica en cuatro niveles:

\begin{enumerate}
  \item \textbf{Representaciones léxicas interpretables.} TF--IDF pondera los
        términos de acuerdo con su frecuencia local y su rareza en el corpus;
        constituye una referencia clásica para recuperación y clasificación de
        texto (Salton \& Buckley, 1988).
  \item \textbf{Modelos lineales supervisados.} Una regresión logística sobre
        unigramas y bigramas permite establecer una línea base rápida,
        reproducible y fácil de inspeccionar. Su desempeño indica cuánto del
        problema puede resolverse mediante asociaciones léxicas.
  \item \textbf{Modelos contextuales.} BERT aprende representaciones
        bidireccionales que pueden adaptarse a clasificación con una capa de
        salida (Devlin et al., 2019); DeBERTa separa información de contenido y
        posición mediante atención desacoplada (He et al., 2021).
  \item \textbf{Adaptación al dominio matemático.} MathBERT fue preentrenado con
        currículo, libros, sílabos y resúmenes matemáticos. En tareas de
        educación matemática, incluida calificación de respuestas abiertas,
        superó a BERT base, lo que evidencia el valor de preservar vocabulario,
        ecuaciones y símbolos especializados (Shen et al., 2021).
\end{enumerate}

La comparación no debe reducirse a escoger el modelo de mayor capacidad. El
benchmark de Otero et al. (2025) encontró que restringir el espacio de
concepciones por tema y revisar errores con educadores produjo mejoras
considerables. Por tanto, contexto, taxonomía y evaluación experta pueden ser
tan importantes como la arquitectura.

\section{Consecuencias para este análisis exploratorio}

La investigación previa conduce a cinco criterios concretos que guían el EDA:

\begin{itemize}
  \item conservar números, fracciones, signos y expresiones matemáticas;
  \item analizar pregunta, opción y explicación como campos distintos;
  \item cuantificar la cola larga y la asociación entre concepción e ítem;
  \item tratar las faltas ortográficas y minúsculas como mezcla de escritura
        auténtica y aumento sintético;
  \item evaluar por separado desempeño en preguntas conocidas y transferencia
        a preguntas nuevas, manteniendo revisión docente.
\end{itemize}

% =============================================================================
\chapter{Marco conceptual de NLP}

\section{Patrones en lenguaje escrito}

El NLP convierte texto en representaciones que un algoritmo puede comparar.
En este reto, una limpieza agresiva sería contraproducente: expresiones como
\codigo{3/9}, \codigo{-13}, \codigo{divide}, \codigo{half} o
\codigo{two negatives} contienen el procedimiento y la posible fuente del
error. Las técnicas relevantes son:

\begin{description}
  \item[Normalización controlada.] Homogeneizar mayúsculas, espacios y LaTeX en
        las opciones sin eliminar números, signos ni barras de fracción.
  \item[Tokenización y n-gramas.] Separar unidades léxicas y conservar
        secuencias cortas como \emph{divided by}, \emph{not shaded} o
        \emph{whole number}, cuyo significado supera al de palabras aisladas.
  \item[TF--IDF.] Ponderar términos frecuentes en una explicación, pero raros
        en el corpus, para obtener una matriz dispersa interpretable
        (Salton \& Buckley, 1988).
  \item[Clasificación lineal.] Combinar TF--IDF con regresión logística como
        línea base fuerte, auditable y eficiente para texto corto.
  \item[Modelos contextuales.] Evaluar posteriormente transformadores como
        BERT o DeBERTa, capaces de representar el contexto de una palabra
        (Devlin et al., 2019; He et al., 2021); su complejidad solo se justifica si
        mejoran una línea base bien validada.
  \item[Clasificación jerárquica.] Separar veracidad, calidad del razonamiento y
        concepción específica, aprovechando la estructura pedagógica de la
        etiqueta en vez de tratar 65 clases como equivalentes.
\end{description}

\section{Por qué el contexto matemático importa}

Una frase no tiene una interpretación universal. ``Se divide por 3'' puede ser
un procedimiento correcto en una pregunta y una regla aplicada indebidamente
en otra. Por ello, el texto debe combinarse con el enunciado y la respuesta.
Los modelos que solo usen \codigo{StudentExplanation} perderán contexto; los
que usen solo \codigo{QuestionId} pueden memorizar asociaciones sin comprender
el razonamiento.

\section{Métrica de la competencia}

El archivo de ejemplo solicita tres etiquetas ordenadas por confianza. Con una
única etiqueta correcta por observación, la precisión promedio a tres
posiciones se reduce al recíproco del rango cuando la respuesta aparece en las
tres primeras posiciones:

\begin{equation}
\operatorname{AP@3}_i =
\begin{cases}
1/r_i, & r_i \leq 3,\\
0, & r_i > 3,
\end{cases}
\qquad
\operatorname{MAP@3}=\frac{1}{n}\sum_{i=1}^{n}\operatorname{AP@3}_i.
\end{equation}

Acertar en la primera, segunda o tercera posición aporta (1), (1/2) o
(1/3), respectivamente. Esta métrica es más útil que la exactitud plana ante
clases raras, aunque no elimina la necesidad de reportar resultados por clase
y por pregunta.

% =============================================================================
\chapter{Descripción de los datos}

\section{Archivos disponibles}

\begin{table}[htbp]
\centering
\small
\caption{Archivos de datos identificados en el proyecto.}
\label{tab:archivos}
\begin{tabularx}{\textwidth}{@{}p{3.7cm}R{2.5cm}R{2.4cm}Y@{}}
\toprule
\textbf{Archivo} & \textbf{N.º obs.} & \textbf{Variables} &
\textbf{Función}\\
\midrule
\codigo{train.csv} & 36,696 & 7 & Exploración y futuro entrenamiento\\
\codigo{test.csv} & 3 & 5 & Ejemplo local del formato de predicción; no permite evaluación real\\
\codigo{sample\_}\linebreak\codigo{submission.csv} & 3 & 2 & Plantilla de tres etiquetas ordenadas por fila\\
Archivo ZIP & -- & -- & Copia comprimida de los tres archivos anteriores\\
\bottomrule
\end{tabularx}
\end{table}

La unidad de análisis no es un estudiante completo ni una calificación final.
Es una respuesta a una pregunta específica. El conjunto no incluye un
identificador de estudiante, curso, centro educativo, edad ni secuencia
temporal; por tanto, no es posible estudiar progreso individual o dependencia
entre respuestas del mismo estudiante.

Además, la unidad debe interpretarse como un \emph{registro de explicación} y
no necesariamente como una producción original independiente. El benchmark
completo combina explicaciones auténticas y paráfrasis sintéticas; la versión
de competencia no proporciona una bandera que permita distinguirlas
(Rittle-Johnson et al., 2025). Esta ausencia condiciona el análisis de
duplicados, ortografía y estilo.

\section{Diccionario de variables}

\small
\begin{longtable}{@{}p{3.5cm}p{3.0cm}p{2.8cm}p{4.7cm}@{}}
\caption{Variables del conjunto de entrenamiento.}
\label{tab:variables}\\
\toprule
\textbf{Variable} & \textbf{Tipo en archivo} & \textbf{Tipo estadístico} &
\textbf{Descripción y rol}\\
\midrule
\endfirsthead
\toprule
\textbf{Variable} & \textbf{Tipo en archivo} & \textbf{Tipo estadístico} &
\textbf{Descripción y rol}\\
\midrule
\endhead
\bottomrule
\endlastfoot
\codigo{row\_id} & Entero & Identificador & Clave única de la observación; no
representa magnitud.\\
\codigo{QuestionId} & Entero & Categórica nominal & Identificador de la
pregunta; codifica contexto y dificultad.\\
\codigo{QuestionText} & Texto & Texto contextual & Enunciado matemático; tres
preguntas incluyen una descripción textual de imagen.\\
\codigo{MC\_Answer} & Texto & Categórica nominal & Opción seleccionada; 96\%
contiene notación LaTeX.\\
\codigo{Student\_}\linebreak\codigo{Explanation} & Texto & Texto libre & Explicación escrita por el
estudiante; predictor lingüístico principal.\\
\codigo{Category} & Texto & Categórica objetivo & Cruce de veracidad
(\codigo{True/False}) y tipo (\codigo{Correct}, \codigo{Misconception} o
\codigo{Neither}).\\
\codigo{Misconception} & Texto & Categórica objetivo & Nombre de la concepción
errónea; solo aplica cuando el tipo es \codigo{Misconception}.\\
\end{longtable}

Los dos enteros son identificadores, no variables cuantitativas. Por eso no
tiene sentido calcular su media como si representaran una escala. Las variables
numéricas útiles se derivan del texto y se presentan en el Capítulo
\ref{ch:eda}.

\section{Estructura de la etiqueta}

\codigo{Category} tiene la forma \codigo{Veracidad\_Tipo}. La veracidad indica
si la opción elegida es correcta; el tipo evalúa la explicación.

\begin{table}[htbp]
\centering
\caption{Interpretación de los seis valores de \codigo{Category}.}
\label{tab:interpretacion-category}
\begin{tabularx}{\textwidth}{@{}p{4.1cm}Y@{}}
\toprule
\textbf{Categoría} & \textbf{Interpretación}\\
\midrule
\codigo{True\_Correct} & Opción correcta y razonamiento correcto.\\
\codigo{True\_Misconception} & Opción correcta, pero explicación basada en una
concepción errónea; posible acierto casual.\\
\codigo{True\_Neither} & Opción correcta, pero explicación insuficiente o no
clasificable como correcta ni como concepción errónea.\\
\codigo{False\_Correct} & Opción incorrecta con procedimiento esencialmente
válido; posible error aritmético final.\\
\codigo{False\_Misconception} & Opción incorrecta y explicación asociada a una
concepción errónea identificable.\\
\codigo{False\_Neither} & Opción incorrecta y explicación insuficiente, irrelevante
o no clasificable dentro de una concepción específica.\\
\bottomrule
\end{tabularx}
\end{table}

% =============================================================================
\chapter{Limpieza y preprocesamiento}

\section{Valores faltantes}

La única columna con ausencias es \codigo{Misconception}: 26,836 registros,
equivalentes a 73.1\% del total. El cruce con el tipo de razonamiento demuestra
una correspondencia perfecta.

\begin{table}[htbp]
\centering
\caption{Disponibilidad de \codigo{Misconception} por tipo de razonamiento.}
\label{tab:nulos}
\begin{tabular}{@{}lrr@{}}
\toprule
\textbf{Tipo} & \textbf{Nulo} & \textbf{Con valor}\\
\midrule
Correct & 15,029 & 0\\
Misconception & 0 & 9,860\\
Neither & 11,807 & 0\\
\midrule
Total & 26,836 & 9,860\\
\bottomrule
\end{tabular}
\end{table}

\begin{decision}[Decisión sobre ausencias]
Los nulos significan ``no aplica'' y se sustituyen por la etiqueta explícita
\codigo{NA}, que coincide con el formato de entrega. No se usa media, moda ni
modelo de imputación: asignar una concepción errónea a casos \codigo{Correct} o
\codigo{Neither} introduciría ruido en casi tres cuartas partes de la base.
\end{decision}

\section{Duplicados y textos repetidos}

El diagnóstico guardado informa:

\begin{itemize}
  \item 0 filas completamente duplicadas cuando se incluye \codigo{row\_id};
  \item 736 filas duplicadas al excluir \codigo{row\_id};
  \item 970 repeticiones adicionales de \codigo{StudentExplanation} (2.6\%);
  \item frases frecuentes como \emph{because it is a whole number} o
        \emph{because two negatives make a positive}.
\end{itemize}

Una repetición puede ser legítima: estudiantes distintos pueden escribir la
misma respuesta breve. Sin embargo, como no existe un identificador de
estudiante, tampoco puede demostrarse que las 736 coincidencias sean
independientes. Se conservan para respetar la unidad original, pero una futura
validación debe evitar que copias idénticas queden a ambos lados de la
partición.

\section{Calidad y convenciones del texto}

\begin{table}[htbp]
\centering
\caption{Indicadores de calidad de \codigo{StudentExplanation}.}
\label{tab:calidad-texto}
\begin{tabularx}{0.86\textwidth}{@{}Y R{3.2cm}@{}}
\toprule
\textbf{Indicador} & \textbf{Proporción}\\
\midrule
Empieza en minúscula & 62.8\%\\
Contiene dígitos & 78.3\%\\
Contiene barra ``/'' & 27.2\%\\
Contiene signo ``='' & 15.9\%\\
Tiene menos de 3 palabras & 0.1\%\\
Contiene LaTeX & 0.0\%\\
\codigo{MC\_Answer} contiene LaTeX & 96.0\%\\
\codigo{QuestionText} contiene descripción de imagen & 23.6\%\\
\bottomrule
\end{tabularx}
\end{table}

El corpus es informal, pero no está vacío ni corrupto. Las faltas ortográficas,
minúsculas y expresiones improvisadas forman parte de la señal observada; no
obstante, una fracción de esos rasgos fue inducida durante el aumento sintético
descrito por Rittle-Johnson et al. (2025). El contenido matemático aparece tanto
en palabras como en números y símbolos.

\section{Decisiones de transformación}

\begin{table}[htbp]
\centering
\caption{Plan de limpieza y justificación.}
\label{tab:limpieza}
\begin{tabularx}{\textwidth}{@{}p{4.0cm}p{3.8cm}Y@{}}
\toprule
\textbf{Situación} & \textbf{Decisión} & \textbf{Justificación}\\
\midrule
Nulos estructurales & Reemplazar por \codigo{NA} & Conserva la semántica y el
formato objetivo.\\
Textos y registros repetidos & Conservar y auditar la partición & Pueden ser
respuestas legítimas, pero no deben causar fuga.\\
LaTeX en la opción & Convertir fracciones a texto plano & Homogeneiza
\codigo{\textbackslash frac\{1\}\{3\}} y \codigo{1/3}.\\
Mayúsculas inconsistentes & Minusculizar la explicación & Reduce variantes sin
cambiar la operación matemática.\\
Dígitos y símbolos & Conservar & Son la señal matemática principal.\\
Ortografía estudiantil & No corregir automáticamente & Una corrección puede
cambiar términos, números o pistas propias de la concepción errónea.\\
Descripciones de imágenes & Conservar como texto & La imagen original no está
disponible; la descripción es el único contexto visual.\\
Espacios repetidos & Colapsar & Evita tokens vacíos sin pérdida semántica.\\
\bottomrule
\end{tabularx}
\end{table}

% =============================================================================
\chapter{Análisis exploratorio}
\label{ch:eda}

\section{Variables categóricas}

\subsection{Distribución de \codigo{Category}}

\begin{table}[htbp]
\centering
\caption{Frecuencia y proporción de las categorías.}
\label{tab:category}
\begin{tabular}{@{}lrr@{}}
\toprule
\textbf{Categoría} & \textbf{Frecuencia} & \textbf{Proporción}\\
\midrule
\codigo{True\_Correct} & 14,802 & 40.34\%\\
\codigo{False\_Misconception} & 9,457 & 25.77\%\\
\codigo{False\_Neither} & 6,542 & 17.83\%\\
\codigo{True\_Neither} & 5,265 & 14.35\%\\
\codigo{True\_Misconception} & 403 & 1.10\%\\
\codigo{False\_Correct} & 227 & 0.62\%\\
\bottomrule
\end{tabular}
\end{table}

\figura{distribucion_category.png}
{Distribución de \codigo{Category} y descomposición entre veracidad y tipo de
razonamiento. Figura extraída de la salida almacenada del notebook.}
{fig:category}

El 55.78\% de las opciones seleccionadas es verdadero y el 44.22\% falso. Al
agrupar por calidad de explicación, 40.96\% es \codigo{Correct}, 32.18\%
\codigo{Neither} y 26.87\% \codigo{Misconception}. La diferencia entre estas
dimensiones demuestra que acertar la opción no equivale a razonar correctamente.

\subsection{Concepciones erróneas y etiqueta compuesta}

Hay 9,860 observaciones con una concepción errónea, distribuidas entre 35
clases. Las más frecuentes se presentan en la Tabla
\ref{tab:misconceptions}.

\begin{table}[htbp]
\centering
\small
\caption{Quince concepciones erróneas más frecuentes.}
\label{tab:misconceptions}
\begin{tabular}{@{}lrr@{}}
\toprule
\textbf{Concepción} & \textbf{Frecuencia} & \textbf{\% entre concepciones}\\
\midrule
Incomplete & 1,454 & 14.75\%\\
Additive & 929 & 9.42\%\\
Duplication & 704 & 7.14\%\\
Subtraction & 620 & 6.29\%\\
Positive & 566 & 5.74\%\\
Wrong\_term & 558 & 5.66\%\\
Irrelevant & 497 & 5.04\%\\
Wrong\_fraction & 418 & 4.24\%\\
Inversion & 414 & 4.20\%\\
Mult & 353 & 3.58\%\\
Denominator-only\_change & 336 & 3.41\%\\
Whole\_numbers\_larger & 329 & 3.34\%\\
Adding\_across & 307 & 3.11\%\\
WNB & 299 & 3.03\%\\
Tacking & 290 & 2.94\%\\
\bottomrule
\end{tabular}
\end{table}

\figura{distribucion_misconceptions.png}
{Concepciones erróneas más frecuentes y cola larga de las 65 etiquetas
compuestas. Las líneas horizontales marcan 50 y 10 observaciones.}
{fig:misconceptions}

La clase mayoritaria \codigo{True\_Correct:NA} tiene 14,802 registros. En el
extremo contrario, la clase compuesta menos frecuente tiene un solo ejemplo.
Treinta y dos clases (49.2\%) tienen menos de 50 observaciones y 24 (36.9\%)
menos de 10. Predecir siempre la clase mayoritaria produciría aproximadamente
40\% de exactitud, por lo que esa métrica aislada sería engañosa.

\begin{advertencia}[Cola larga]
Las clases con menos de 10 casos no pueden dividirse de manera estable entre
entrenamiento y validación. Cualquier resultado agregado debe acompañarse de
soporte por clase y cobertura; sobremuestrear no crea diversidad lingüística
nueva y puede aumentar la memorización.
\end{advertencia}

\section{Variables cuantitativas derivadas del texto}

Debido a que la base no contiene medidas numéricas nativas, se construyeron
características que describen forma y contenido de la explicación.

\begin{table}[htbp]
\centering
\caption{Variables cuantitativas derivadas.}
\label{tab:derivadas}
\begin{tabularx}{\textwidth}{@{}p{3.2cm}Y@{}}
\toprule
\textbf{Variable} & \textbf{Definición}\\
\midrule
\codigo{n\_chars} & Número de caracteres de la explicación normalizada.\\
\codigo{n\_words} & Número de palabras o tokens separados por espacios.\\
\codigo{n\_unique\_w} & Número de palabras distintas.\\
\codigo{avg\_token\_len} & Longitud media de los tokens separados por espacios;
incluye signos adheridos al token.\\
\codigo{lex\_diversity} & Cociente entre palabras distintas y palabras totales.\\
\codigo{n\_digits} & Cantidad de dígitos.\\
\codigo{digit\_ratio} & Proporción de caracteres que son dígitos.\\
\codigo{n\_math\_sym} & Conteo de símbolos como $+$, $-$, $*$, $/$, $=$,
$<$, $>$ y \%.\\
\codigo{n\_sentences} & Aproximación al número de segmentos u oraciones.\\
\bottomrule
\end{tabularx}
\end{table}

\begin{table}[htbp]
\centering
\scriptsize
\caption{Resumen estadístico de las variables derivadas.}
\label{tab:descriptivas}
\resizebox{\textwidth}{!}{%
\begin{tabular}{@{}lrrrrrrrrrr@{}}
\toprule
\textbf{Variable} & \textbf{Media} & \textbf{DE} & \textbf{Mín.} &
\textbf{Q1} & \textbf{Mediana} & \textbf{Q3} & \textbf{Máx.} &
\textbf{IQR} & \textbf{CV} & \textbf{Asimetría}\\
\midrule
n\_chars & 69.78 & 38.44 & 1.00 & 43.00 & 60.00 & 86.00 & 586.00 & 43.00 & 0.55 & 1.89\\
n\_words & 15.36 & 8.12 & 1.00 & 10.00 & 13.00 & 19.00 & 177.00 & 9.00 & 0.53 & 1.87\\
n\_unique\_w & 13.51 & 6.22 & 1.00 & 9.00 & 12.00 & 17.00 & 125.00 & 8.00 & 0.46 & 1.56\\
avg\_token\_len & 3.64 & 0.68 & 1.00 & 3.20 & 3.62 & 4.00 & 14.83 & 0.80 & 0.19 & 0.84\\
lex\_diversity & 0.91 & 0.09 & 0.39 & 0.85 & 0.92 & 1.00 & 1.00 & 0.15 & 0.10 & -0.58\\
n\_digits & 4.72 & 4.61 & 0.00 & 1.00 & 4.00 & 7.00 & 101.00 & 6.00 & 0.98 & 1.63\\
digit\_ratio & 0.07 & 0.07 & 0.00 & 0.02 & 0.06 & 0.11 & 0.78 & 0.09 & 0.96 & 1.34\\
n\_math\_sym & 1.01 & 1.66 & 0.00 & 0.00 & 0.00 & 2.00 & 50.00 & 2.00 & 1.64 & 2.92\\
n\_sentences & 1.11 & 0.41 & 1.00 & 1.00 & 1.00 & 1.00 & 11.00 & 0.00 & 0.37 & 5.23\\
\bottomrule
\end{tabular}}
\end{table}

\figura{distribuciones_variables_texto.png}
{Histogramas de seis variables cuantitativas derivadas. Las explicaciones son
cortas y las medidas de longitud presentan cola derecha.}
{fig:distribuciones-texto}

La mediana es de 13 palabras y 60 caracteres, frente a máximos de 177 palabras
y 586 caracteres. La diferencia entre media y mediana, junto con asimetrías de
1.89 y 1.87, confirma cola derecha. La mediana de símbolos matemáticos es cero,
aunque una minoría concentra muchos. La aproximación revisada de
\codigo{n\_sentences} separa únicamente puntuación terminal seguida de espacio o
fin de texto: su máximo es 11 y su mediana es 1. Sigue siendo una variable
heurística, por lo que debe interpretarse con cautela.

\section{Longitud y tipo de razonamiento}

\begin{table}[htbp]
\centering
\caption{Longitud de explicación por categoría.}
\label{tab:longitud-category}
\begin{tabular}{@{}lrrr@{}}
\toprule
\textbf{Categoría} & \textbf{Media} & \textbf{Mediana} & \textbf{DE}\\
\midrule
False\_Correct & 18.55 & 17 & 8.23\\
True\_Correct & 17.62 & 16 & 9.03\\
True\_Misconception & 16.40 & 15 & 8.90\\
False\_Misconception & 14.80 & 13 & 6.86\\
True\_Neither & 13.04 & 11 & 7.38\\
False\_Neither & 12.74 & 11 & 6.53\\
\bottomrule
\end{tabular}
\end{table}

\figura{longitud_por_categoria.png}
{Diagramas de caja y violín para la longitud de la explicación. Las medianas
por tipo son 16 palabras para \codigo{Correct}, 13 para
\codigo{Misconception} y 11 para \codigo{Neither}.}
{fig:longitud-category}

Los razonamientos correctos tienden a ser más extensos, posiblemente porque
describen un procedimiento. Las explicaciones \codigo{Neither} son más cortas
y con frecuencia no aportan evidencia evaluable. No obstante, las
distribuciones se solapan ampliamente: la longitud es una característica
auxiliar, no una regla de clasificación.

\section{Valores atípicos}

El criterio de rango intercuartílico $1.5\times IQR$ identificó los siguientes
extremos:

\begin{table}[htbp]
\centering
\small
\caption{Atípicos según el criterio IQR.}
\label{tab:outliers}
\begin{tabular}{@{}lrrrr@{}}
\toprule
\textbf{Variable} & \textbf{Límite inferior} & \textbf{Límite superior} &
\textbf{Atípicos} & \textbf{Proporción}\\
\midrule
n\_chars & -21.50 & 150.50 & 1,556 & 4.24\%\\
n\_words & -3.50 & 32.50 & 1,499 & 4.08\%\\
n\_unique\_w & -3.00 & 29.00 & 815 & 2.22\%\\
avg\_token\_len & 2.00 & 5.20 & 851 & 2.32\%\\
lex\_diversity & 0.62 & 1.23 & 12 & 0.03\%\\
n\_digits & -8.00 & 16.00 & 742 & 2.02\%\\
digit\_ratio & -0.12 & 0.25 & 831 & 2.26\%\\
n\_math\_sym & -3.00 & 5.00 & 882 & 2.40\%\\
n\_sentences & 1.00 & 1.00 & 3,097 & 8.44\%\\
\bottomrule
\end{tabular}
\end{table}

Una explicación larga puede ser un procedimiento detallado y una explicación
muy corta puede indicar respuesta mínima o ausencia de razonamiento. Hay 30
observaciones con dos palabras o menos (0.08\%). También existe al menos una
explicación de 177 palabras con contenido que parece una secuencia extensa de
símbolos y asociaciones, por lo que no todos los extremos deben describirse
automáticamente como alta calidad. La inspección contextual es necesaria.

\begin{decision}[Tratamiento de atípicos]
No se eliminan solo por superar el umbral IQR. En texto humano, la longitud
extrema puede ser señal predictiva. Para modelar conviene transformar
\codigo{n\_words} con \(\log(1+x)\), limitar la influencia numérica si fuera
necesario y auditar manualmente los casos más extremos.
\end{decision}

\section{Correlaciones}

Se utiliza Spearman porque las variables son asimétricas y las relaciones
pueden ser monótonas sin ser lineales.

\figura{correlaciones_variables_texto.png}
{Matriz de correlaciones de Spearman y relación entre palabras totales y
palabras únicas.}
{fig:correlaciones}

Los pares más correlacionados son \codigo{n\_words} con
\codigo{n\_unique\_w} (\(\rho=0.981\)), \codigo{n\_chars} con
\codigo{n\_words} (\(\rho=0.951\)) y \codigo{n\_chars} con
\codigo{n\_unique\_w} (\(\rho=0.945\)). \codigo{n\_digits} y
\codigo{digit\_ratio} alcanzan \(\rho=0.894\). La diversidad léxica tiene
correlación negativa moderada con el número de palabras
(\(\rho\approx-0.59\)), en parte porque repetir alguna palabra se vuelve más
probable conforme crece el texto.

\begin{hallazgo}[Redundancia]
Incluir caracteres, palabras y palabras únicas simultáneamente aportaría
información muy similar. Una alternativa parsimoniosa es conservar
\codigo{n\_words}, una medida normalizada como \codigo{digit\_ratio} y
banderas matemáticas específicas.
\end{hallazgo}

\section{Cruces entre pregunta, respuesta y etiqueta}

\subsection{Dificultad heterogénea}

Las 15 preguntas contienen entre 673 y 4,857 observaciones. La tasa de acierto
varía entre 35.7\% y 79.6\%, con media no ponderada de 56.7\%. Por ello,
\codigo{QuestionId} contiene información contextual y de dificultad; no debe
interpretarse como una cantidad continua.

\figura{cruces_por_pregunta.png}
{Composición de las categorías y tasa de acierto para cada pregunta. La línea
vertical discontinua marca 50\% de respuestas correctas.}
{fig:cruces-pregunta}

La variabilidad se observa tanto en el acierto como en la proporción de
\codigo{Neither} y \codigo{Misconception}. Una pregunta difícil no solo genera
más respuestas falsas; puede inducir patrones de error distintos.

\subsection{Concepciones ligadas a la pregunta}

De las 35 concepciones erróneas, 32 aparecen en una sola pregunta, dos en dos
preguntas y una en tres preguntas. Solo \codigo{Incomplete},
\codigo{Irrelevant} y \codigo{WNB} se describen como transversales. El 91\% de
las concepciones es, por tanto, exclusivo de un contexto local.

Esto crea dos escenarios de validación:

\begin{enumerate}
  \item Una partición aleatoria conserva las mismas preguntas a ambos lados y
        estima desempeño en preguntas conocidas, pero puede premiar la
        memorización pregunta--etiqueta.
  \item Una partición agrupada por \codigo{QuestionId} evalúa preguntas nuevas,
        pero deja fuera del entrenamiento casi todas las concepciones que
        aparecen en validación; muchas resultan imposibles de predecir por
        construcción.
\end{enumerate}

Con solo 15 preguntas, ninguno de los dos esquemas por sí solo describe toda la
capacidad del sistema.

\subsection{Veracidad casi determinista}

Existen 60 combinaciones de pregunta y opción. En 58, la veracidad es única.
Las dos excepciones pertenecen a la pregunta 31778:

\begin{table}[htbp]
\centering
\caption{Pares ambiguos de pregunta y respuesta.}
\label{tab:ambiguos}
\begin{tabular}{@{}lrrr@{}}
\toprule
\textbf{Pregunta} & \textbf{Opción} & \textbf{True} & \textbf{False}\\
\midrule
31778 & 6 & 1,749 & 6\\
31778 & 9 & 12 & 138\\
\bottomrule
\end{tabular}
\end{table}

Las minorías pueden representar ruido de etiquetado o una ambigüedad de
anotación. En cualquier caso, la veracidad se resuelve casi por completo con
una tabla de consulta; el NLP debe concentrarse en la calidad y naturaleza del
razonamiento.

\section{Contenido léxico}

Una representación TF--IDF con unigramas y bigramas, vocabulario máximo de
3,000 términos y frecuencia mínima de cinco produce una matriz de
36,696 por 3,000 con densidad de 0.332\%. La dispersión es esperable en textos
de mediana 13 palabras.

Los términos almacenados como característicos incluyen \emph{divided},
\emph{times} y expresiones numéricas para \emph{Correct}; \emph{not shaded},
\emph{positive} y \emph{add} para \codigo{Misconception}; y \emph{worked out},
\emph{guess}, \emph{sure}, \emph{don't know} o \emph{explain} para
\codigo{Neither}.

\figura[0.91\textwidth]{marcadores_lexicos.png}
{Composición del tipo de razonamiento condicionada a palabras seleccionadas.
La línea discontinua representa la tasa base de razonamientos correctos.}
{fig:marcadores-lexicos}

\emph{Divide}, \emph{multiply}, \emph{fraction}, \emph{times} y \emph{equal}
se asocian con una proporción de razonamiento correcto superior a la base;
\emph{guess} se concentra claramente en \codigo{Neither}. Sin embargo, algunas
palabras también identifican la pregunta concreta. El poder aparente del
vocabulario debe comprobarse con una validación que separe contexto de
generalización.

% =============================================================================
\chapter{Implicaciones para validación y modelado}

\section{Amenazas de fuga de información}

La repetición de preguntas, opciones y algunas explicaciones puede inflar una
partición aleatoria. Una predicción correcta no demuestra comprensión si el
modelo aprendió que cierta opción de cierta pregunta suele acompañarse de una
etiqueta específica. Deben distinguirse tres fuentes de señal:

\begin{enumerate}
  \item \textbf{Regla del ítem:} \codigo{QuestionId + MC\_Answer} determina la
        veracidad en 96.7\% de los pares distintos.
  \item \textbf{Inventario local:} la pregunta restringe las concepciones
        plausibles.
  \item \textbf{Lenguaje del razonamiento:} palabras, operaciones y expresiones
        distinguen cómo llegó el estudiante a la respuesta.
\end{enumerate}

Un sistema útil puede explotar las tres, pero su evaluación debe indicar qué
tipo de generalización mide.

\section{Propuesta de validación doble}

\begin{table}[htbp]
\centering
\caption{Esquemas complementarios de validación.}
\label{tab:validacion}
\begin{tabularx}{\textwidth}{@{}p{3.4cm}p{3.6cm}Y@{}}
\toprule
\textbf{Esquema} & \textbf{Qué estima} & \textbf{Limitación}\\
\midrule
\codigo{StratifiedKFold} sobre etiqueta compuesta & Desempeño en preguntas y
clases ya conocidas & Puede beneficiarse de memorizar pregunta, opción o texto
repetido.\\
\codigo{GroupKFold} por \codigo{QuestionId} & Transferencia a preguntas nuevas
& El 91\% de concepciones del pliegue puede no existir en entrenamiento.\\
Agrupación adicional de duplicados & Robustez ante textos idénticos & Reduce
ligeramente el tamaño útil y requiere definir una clave de grupo.\\
\bottomrule
\end{tabularx}
\end{table}

La distancia entre los dos primeros resultados es un hallazgo, no un error:
cuantifica cuánto depende el sistema de preguntas conocidas.

\section{Arquitectura jerárquica propuesta}

\begin{enumerate}
  \item \textbf{Etapa 1: veracidad.} Usar el par
        \codigo{(QuestionId, MC\_Answer)} y tratar explícitamente los dos pares
        ambiguos.
  \item \textbf{Etapa 2: tipo.} Clasificar la explicación como
        \codigo{Correct}, \codigo{Misconception} o \codigo{Neither}, combinando
        texto, pregunta y opción.
  \item \textbf{Etapa 3: concepción.} Si el tipo es
        \codigo{Misconception}, ordenar las concepciones plausibles. Para
        preguntas conocidas puede restringirse el inventario; para preguntas
        nuevas debe mantenerse una ruta abierta o taxonomía más general.
  \item \textbf{Composición.} Convertir las probabilidades en tres etiquetas
        \codigo{Category:Misconception} ordenadas para MAP@3.
\end{enumerate}

\section{Línea base e ingeniería de características}

Antes de un transformador, se recomienda establecer TF--IDF de unigramas y
bigramas con regresión logística. La línea base debe integrar texto de la
pregunta, opción normalizada y explicación, distinguiendo cada campo con
prefijos. Como variables auxiliares pueden incluirse:

\begin{itemize}
  \item \(\log(1+\codigo{n\_words})\);
  \item \codigo{digit\_ratio};
  \item presencia de fracción y operadores;
  \item dificultad estimada de la pregunta, calculada solo con el pliegue de
        entrenamiento;
  \item longitud relativa frente a otras explicaciones de la misma pregunta.
\end{itemize}

No conviene incluir simultáneamente todas las medidas de longitud por su alta
correlación. Para la cola larga pueden evaluarse pesos de clase, agrupaciones
pedagógicamente justificadas o recuperación de candidatos; un simple
sobremuestreo duplica ejemplos, pero no añade variación lingüística.

% =============================================================================
\chapter{Hallazgos, conclusiones y recomendaciones}

\section{Síntesis de hallazgos}

\begin{longtable}{@{}C{0.8cm}p{4.4cm}p{4.1cm}p{4.2cm}@{}}
\caption{Hallazgos principales y consecuencias.}
\label{tab:hallazgos}\\
\toprule
\textbf{N.º} & \textbf{Hallazgo} & \textbf{Evidencia} & \textbf{Consecuencia}\\
\midrule
\endfirsthead
\toprule
\textbf{N.º} & \textbf{Hallazgo} & \textbf{Evidencia} & \textbf{Consecuencia}\\
\midrule
\endhead
\bottomrule
\endlastfoot
1 & Los nulos son estructurales & Correspondencia perfecta con tipo
\codigo{Misconception} & Codificar \codigo{NA}; no imputar.\\
2 & Hay repeticiones que requieren cautela & 736 duplicados sin
\codigo{row\_id}; 970 textos repetidos & Conservar, pero agrupar copias en la
validación.\\
3 & El texto es informal y matemático & 62.8\% inicia en minúscula; 78.3\% tiene
dígitos & No eliminar números ni símbolos.\\
4 & \codigo{Category} cruza dos dimensiones & Veracidad por tipo de razonamiento
& Favorece clasificación jerárquica.\\
5 & Existe desbalance extremo & 32/65 clases con menos de 50 casos; 24 con
menos de 10 & Usar MAP@3, soporte por clase y estrategias de cola larga.\\
6 & La longitud discrimina débilmente & Medianas 16, 13 y 11 palabras & Usarla
como variable auxiliar, no como regla.\\
7 & Los atípicos pueden ser informativos & 4.08\% de atípicos en palabras por
IQR & Auditar y transformar; no eliminar automáticamente.\\
8 & Las longitudes son redundantes & Correlaciones de Spearman superiores a
0.94 & Elegir representación parsimoniosa.\\
9 & La dificultad cambia por pregunta & Acierto entre 35.7\% y 79.6\% & Tratar
\codigo{QuestionId} como contexto categórico.\\
10 & Las concepciones están ligadas al ítem & 32/35 aparecen en una pregunta &
La validación debe separar preguntas conocidas y nuevas.\\
11 & La veracidad es casi determinista & 58/60 pares pregunta--opción sin
ambigüedad & Separar veracidad del análisis del razonamiento.\\
12 & El vocabulario contiene señal & TF--IDF de 3,000 términos y asociaciones
interpretables & Línea base clásica viable antes de transformadores.\\
\end{longtable}

\section{Conclusiones}

\begin{enumerate}
  \item El conjunto tiene una estructura clara y no presenta ausencias reales
        en las variables obligatorias. La principal tarea de limpieza es
        interpretar correctamente los nulos, normalizar convenciones y
        conservar la notación matemática.
  \item La procedencia es mixta: el benchmark completo incluye 27\% de
        explicaciones aumentadas y los archivos no permiten distinguirlas. Por
        ello, la calidad textual no puede atribuirse exclusivamente a estudiantes.
  \item La etiqueta no representa un único fenómeno. Separa la corrección de la
        opción y la calidad del razonamiento, lo que permite detectar aciertos
        casuales y errores finales de cálculo.
  \item El desbalance es suficientemente severo para invalidar una lectura
        basada solo en exactitud. La cola larga debe analizarse con soporte,
        métricas de ranking y resultados desagregados.
  \item La explicación contiene señal útil: longitud, símbolos y vocabulario
        cambian entre tipos de razonamiento. Sin embargo, ninguna característica
        superficial basta por sí sola y los grupos se solapan.
  \item La amenaza principal a la validez es la dependencia entre pregunta y
        concepción. Un buen resultado aleatorio puede reflejar memorización; un
        resultado por pregunta puede penalizar clases nunca observadas.
  \item La veracidad puede resolverse casi por completo sin NLP, mientras que
        interpretar el razonamiento y asignar la concepción específica concentra
        la dificultad científica.
  \item Una línea base TF--IDF más regresión logística, evaluada con dos esquemas
        de validación, ofrece el punto de comparación necesario antes de modelos
        contextuales más costosos.
\end{enumerate}

\section{Recomendaciones y próximos pasos}

\begin{enumerate}
  \item Corregir la ruta de datos del notebook y ejecutar posteriormente una
        prueba de reproducibilidad en un ambiente limpio, registrando versiones
        de dependencias y semilla aleatoria.
  \item Revisar los 736 registros idénticos sin identificador y construir grupos
        que impidan su separación entre entrenamiento y validación.
  \item Reportar \codigo{StratifiedKFold} y \codigo{GroupKFold} por pregunta,
        además de MAP@3 global, soporte por clase y métricas por pregunta.
  \item Construir el modelo jerárquico en etapas y comparar contra un
        clasificador plano de 65 clases.
  \item Establecer TF--IDF con regresión logística antes de evaluar DeBERTa o un
        modelo especializado en lenguaje matemático.
  \item Analizar manualmente errores por categoría, especialmente
        \codigo{True\_}\allowbreak\codigo{Misconception},
        \codigo{False\_}\allowbreak\codigo{Correct} y clases con menos
        de 10 ejemplos.
  \item Mantener supervisión docente: la predicción debe orientar revisión y
        retroalimentación, no convertirse en un diagnóstico pedagógico
        incuestionable.
\end{enumerate}

\section{Limitaciones}

\begin{itemize}
  \item Solo hay 15 preguntas, por lo que la variedad de conceptos y contextos
        es limitada.
  \item El archivo de prueba local contiene tres filas y no permite evaluar
        generalización externa.
  \item No existe identificador de estudiante; no se puede medir dependencia,
        progreso ni repetición individual.
  \item No existe una variable que identifique las paráfrasis sintéticas; no es
        posible comparar su distribución, dificultad o efecto por clase.
  \item Tres preguntas dependen de una descripción textual de imagen, no de la
        imagen original.
  \item Las etiquetas raras tienen soporte insuficiente para estimaciones
        estables.
  \item Los términos TF--IDF pueden reflejar el tema de la pregunta además del
        tipo de razonamiento.
  \item El análisis de este informe usa salidas almacenadas y no sustituye una
        futura auditoría de ejecución completa.
\end{itemize}

\begin{decision}[Cierre]
La evidencia cumple el propósito exploratorio: describe variables y
observaciones, justifica la limpieza, estudia categorías y medidas derivadas,
analiza correlaciones y atípicos, cruza los factores esenciales y convierte los
hallazgos en un plan verificable. El siguiente avance no debe comenzar por el
modelo más complejo, sino por una validación honesta y una línea base
interpretable.
\end{decision}

% =============================================================================
\appendix

\chapter{Matriz de cumplimiento de la rúbrica}

\begin{longtable}{@{}p{4.0cm}p{9.8cm}@{}}
\caption{Correspondencia entre los criterios de evaluación y el informe.}
\label{tab:cumplimiento}\\
\toprule
\textbf{Criterio} & \textbf{Evidencia en el documento}\\
\midrule
\endfirsthead
\toprule
\textbf{Criterio} & \textbf{Evidencia en el documento}\\
\midrule
\endhead
\bottomrule
\endlastfoot
Situación problemática & Capítulo 1: contexto educativo, relevancia,
diferencia entre resultado y razonamiento, y valor de detectar aciertos
casuales.\\
Problema científico & Pregunta específica y formulación formal de predictores
y etiqueta en el Capítulo 1.\\
Objetivos & Un objetivo general y cuatro objetivos específicos medibles y
alcanzables.\\
Investigación previa & Capítulos 2 y 3: competencia, procedencia del benchmark,
anotación, aumento sintético y técnicas de NLP, con citas y referencias APA 7.\\
Descripción de los datos & Capítulo 4: archivos, unidad de análisis, 36,696
observaciones, siete variables, diccionario e interpretación de etiquetas.\\
Limpieza y preprocesamiento & Capítulo 5: ausencias estructurales, duplicados,
calidad del texto y tabla de decisiones justificadas.\\
Estadística descriptiva & Tablas \ref{tab:category},
\ref{tab:misconceptions}, \ref{tab:descriptivas} y
\ref{tab:longitud-category}.\\
Gráficos exploratorios & Figuras \ref{fig:category}--\ref{fig:marcadores-lexicos}:
barras, mapa de calor, histogramas, cajas, violines, dispersión y correlación.\\
Correlaciones & Sección 6.5 y Figura \ref{fig:correlaciones}, con elección
justificada de Spearman.\\
Outliers & Sección 6.4 y Tabla \ref{tab:outliers}; explicación contextual y
decisión de conservación.\\
Variables categóricas & Frecuencias, proporciones, desbalance y cruces por
pregunta en el Capítulo 6.\\
Hallazgos y conclusiones & Capítulo 8: síntesis de doce hallazgos, conclusiones
actualizadas y próximos pasos.\\
Implicaciones & Capítulo 7: fuga, validación doble, arquitectura jerárquica e
ingeniería de características.\\
Informe sin código & El cuerpo explica procedimientos y resultados sin incluir
celdas de programación.\\
\end{longtable}

\chapter{Notas de consistencia y reproducibilidad}

Al convertir el notebook en informe se comparó la narrativa con las salidas
almacenadas. Se adoptaron los siguientes valores verificables:

\begin{itemize}
  \item asimetría de \codigo{n\_chars}=1.89 y
        \codigo{n\_words}=1.87, no valores aproximados superiores a 2;
  \item correlación de Spearman entre \codigo{lex\_diversity} y
        \codigo{n\_words} cercana a -0.59, no -0.85;
  \item densidad TF--IDF de 0.332\%, según la salida de la matriz;
  \item tasa de acierto por pregunta entre 35.7\% y 79.6\%, no un rango de
        20\% a 85\%;
  \item existencia de 736 duplicados al ignorar \codigo{row\_id}, aunque no
        existan filas completas idénticas.
\end{itemize}

Estas correcciones no alteran las conclusiones centrales, pero mejoran la
trazabilidad. Para una entrega definitiva se recomienda ejecutar el notebook de
principio a fin después de corregir la ruta de datos y comprobar que texto,
tablas y figuras se actualicen de forma automática.

\backmatter

\chapter*{Referencias}
\addcontentsline{toc}{chapter}{Referencias}

\begin{list}{}{%
  \setlength{\leftmargin}{1.27cm}
  \setlength{\itemindent}{-1.27cm}
  \setlength{\itemsep}{0.85em}
  \setlength{\parsep}{0pt}}

\item Devlin, J., Chang, M.-W., Lee, K., \& Toutanova, K. (2019). BERT:
Pre-training of deep bidirectional transformers for language understanding. En
\emph{Proceedings of the 2019 Conference of the North American Chapter of the
Association for Computational Linguistics: Human Language Technologies}
(pp. 4171--4186). Association for Computational Linguistics.
\url{https://doi.org/10.18653/v1/N19-1423}

\item García Pérez, L. (2026). \emph{Guía del Proyecto 2: Análisis
Exploratorio} [Guía de proyecto]. Universidad del Valle de Guatemala.

\item Grupo 1, Sección 10. (2026). \emph{Proyecto 2: MAP -- Charting Student
Math Misunderstandings} [Repositorio de GitHub]. GitHub.
\url{https://github.com/DanielBarillasM/Proyecto-2_Grupo-1_DS_Sec-10}

\item He, P., Liu, X., Gao, J., \& Chen, W. (2021). DeBERTa:
Decoding-enhanced BERT with disentangled attention. En \emph{International
Conference on Learning Representations}.
\url{https://openreview.net/forum?id=XPZIaotutsD}

\item Imran, M., \& Bulathwela, S. (2026). \emph{The correct answer trap:
Pedagogically-grounded detection and feedback for hidden misconceptions}
[Trabajo de taller]. arXiv. \url{https://arxiv.org/abs/2606.23205}

\item King, J., Smith, K., Burleigh, L., Crossley, S., Demkin, M., \& Reade,
W. (2025). \emph{MAP -- Charting Student Math Misunderstandings} [Competencia
de ciencia de datos]. Kaggle.
\url{https://www.kaggle.com/competitions/map-charting-student-math-misunderstandings}

\item Otero, N., Druga, S., \& Lan, A. (2025). A benchmark for math
misconceptions: Bridging gaps in middle school algebra with AI-supported
instruction. \emph{Discover Education, 4}, artículo 277.
\url{https://doi.org/10.1007/s44217-025-00742-w}

\item Rittle-Johnson, B., Adler, R., Durkin, K., Burleigh, L., King, J., \&
Crossley, S. (2025). Detecting math misconceptions: An AI benchmark dataset.
En \emph{Proceedings of the Artificial Intelligence in Measurement and
Education Conference (AIME-Con): Works in Progress} (pp. 20--24). National
Council on Measurement in Education.
\url{https://aclanthology.org/2025.aimecon-wip.3/}

\item Salton, G., \& Buckley, C. (1988). Term-weighting approaches in automatic
text retrieval. \emph{Information Processing \& Management, 24}(5), 513--523.
\url{https://doi.org/10.1016/0306-4573(88)90021-0}

\item Shen, J. T., Yamashita, M., Prihar, E., Heffernan, N., Wu, X., Graff, B.,
\& Lee, D. (2021). \emph{MathBERT: A pre-trained language model for general NLP
tasks in mathematics education} [Preprint]. arXiv.
\url{https://doi.org/10.48550/arXiv.2106.07340}

\item The Learning Agency. (2025, 12 de noviembre). \emph{Case study: Math
misconceptions competition}.
\href{https://the-learning-agency.com/the-cutting-ed/article/case-study-math-misconceptions-competition/}{the-learning-agency.com/the-cutting-ed/}

\end{list}

\vfill
\begin{center}
\colorbox{AzulProfundo}{%
  \begin{minipage}{0.72\textwidth}
  \centering\vspace{2.5mm}\color{white}\sffamily
  \textbf{Fin del informe}\vspace{2.5mm}
  \end{minipage}}
\end{center}

\end{document}
